\documentclass[11pt]{article}
\usepackage[a4paper,margin=2.2cm]{geometry}
\usepackage{amsmath,amssymb,amsfonts}
\usepackage{fontspec}
\usepackage{xcolor}
\usepackage{tcolorbox}
\usepackage{booktabs}
\usepackage{enumitem}
\usepackage{array}
\tcbuselibrary{skins,breakable}
\setmainfont{PingFang SC}
\XeTeXlinebreaklocale "zh"
\XeTeXlinebreakskip = 0pt plus 1pt
\definecolor{okgreen}{HTML}{1B7F3B}
\definecolor{warnred}{HTML}{B3261E}
\definecolor{aiblue}{HTML}{1F4E79}
\definecolor{rulegray}{HTML}{7A7A7A}
\newcommand{\okmark}{\textcolor{okgreen}{\textbf{[已核验]}}}
\newcommand{\warnmark}{\textcolor{warnred}{\textbf{[未通过]}}}
\newtcolorbox{aibox}[1]{colback=aiblue!4,colframe=aiblue!55,
  boxrule=0.5pt,arc=2pt,breakable,title=#1,fonttitle=\bfseries\small}
\newtcolorbox{veribox}[1]{colback=okgreen!4,colframe=okgreen!45,
  boxrule=0.5pt,arc=2pt,breakable,title=#1,fonttitle=\bfseries\small}

\begin{document}
\begin{center}
{\LARGE\bfseries 线性代数与大学物理 作业自动求解报告}\\[6pt]
{\small 课程：线性代数与大学物理 \quad\textbar\quad 学生：Zhouruoying \quad\textbar\quad 题数：13 \quad\textbar\quad 核验通过：11/13 \quad\textbar\quad 生成时间：2026-10-08 16:32:35}
\end{center}
\vspace{6pt}
\hrule
\vspace{10pt}
\section*{一、题目与解答}

\subsection*{题 P01}
\noindent\textbf{题目.}\quad 1. 求解下列线性方程组
求解下列三元线性方程组：

\[\begin{cases}
3x + 2y - z = 4 \\
x - y + 2z = 5 \\
2x + 3y + z = 11
\end{cases}\]

要求给出唯一解，并说明消元路径。\par
\vspace{4pt}\noindent\textbf{解答.}\quad
\begin{enumerate}[leftmargin=1.6em,itemsep=2pt]
  \item 写成增广矩阵 $\left[\,A \mid \mathbf{b}\,\right]$ 并对左侧作初等行变换：
  \item 化为行阶梯形后消元，得到唯一解。
  \item 回代校验：将解代回每一个原方程，左端与右端残差均为 $0$。
\end{enumerate}
\[ \boxed{x = 1,\quad y = 2,\quad z = 3} \]
\noindent{\small 方法：高斯消元 / \texttt{sympy.solve}}\par
\begin{veribox}{独立核验 \okmark}substitute：所有方程回代残差为 0\end{veribox}
\begin{aibox}{AI 讲解（本地 qwen3:8b，仅讲解不产数值）}通过消元法将方程组化为阶梯形，进而求得唯一解。\end{aibox}
\vspace{10pt}
\subsection*{题 P02}
\noindent\textbf{题目.}\quad 2. 求矩阵的特征值与特征向量
求矩阵 $A = \begin{pmatrix} 4 & 1 \\ 2 & 3 \end{pmatrix}$ 的特征值与对应的特征向量。\par
\vspace{4pt}\noindent\textbf{解答.}\quad
\begin{enumerate}[leftmargin=1.6em,itemsep=2pt]
  \item 解特征方程 $\det(A-\lambda I)=0$：
  \item $\lambda = 5$（代数重数 1），$\lambda = 2$（代数重数 1）。
  \item 把 $\lambda$ 代回 $(A-\lambda I)\mathbf{x}=0$，取基础解系即为特征向量。
\end{enumerate}
\[ \boxed{\lambda_{1}=2,\quad \lambda_{2}=5;\quad \text{对应特征向量见下}} \]
\noindent{\small 方法：特征方程 + 齐次方程组基础解系}\par
\begin{veribox}{独立核验 \okmark}eigen\_residual：2 组特征对满足 Av=λv\end{veribox}
\begin{aibox}{AI 讲解（本地 qwen3:8b，仅讲解不产数值）}通过求解特征方程得到特征值，再代入原矩阵求对应特征向量。\end{aibox}
\vspace{10pt}
\subsection*{题 P03}
\noindent\textbf{题目.}\quad 3. 求矩阵的秩
求矩阵 $A = \begin{pmatrix} 1 & 2 & 3 \\ 2 & 4 & 6 \\ 3 & 6 & 10 \end{pmatrix}$ 的秩，并说明行向量组是否线性相关。\par
\vspace{4pt}\noindent\textbf{解答.}\quad
\begin{enumerate}[leftmargin=1.6em,itemsep=2pt]
  \item 对 $A$ 作初等行变换化为行阶梯形：非零行数即为矩阵的秩。
  \item 行阶梯形中非零行共 $2$ 行，故 $r(A)=2$。
  \item 其中 $1$ 行被完全消去，说明行向量组线性相关。
\end{enumerate}
\[ \boxed{r(A)=2} \]
\noindent{\small 方法：初等行变换化行阶梯形}\par
\begin{veribox}{独立核验 \okmark}rank\_alt：sympy 秩=2, 非零子式最高阶=2, numpy SVD 秩=2\end{veribox}
\begin{aibox}{AI 讲解（本地 qwen3:8b，仅讲解不产数值）}通过行变换化简矩阵，确定其秩并判断行向量组是否线性相关。\end{aibox}
\vspace{10pt}
\subsection*{题 P04}
\noindent\textbf{题目.}\quad 4. 计算行列式的值
计算下列三阶行列式的值：$\begin{vmatrix} 2 & 1 & 3 \\ 4 & 5 & 6 \\ 7 & 8 & 10 \end{vmatrix}$。\par
\vspace{4pt}\noindent\textbf{解答.}\quad
\begin{enumerate}[leftmargin=1.6em,itemsep=2pt]
  \item $n=3$ 阶行列式，按第一行（或含零最多的行）展开降阶。
  \item 逐项计算代数余子式并代数求和。
  \item 结果见答案框（已通过两条独立算法复核）。
\end{enumerate}
\[ \boxed{\det A=-3} \]
\noindent{\small 方法：Laplace 展开 + 独立算法复核}\par
\begin{veribox}{独立核验 \okmark}det\_alt：Bareiss=-3, Berkowitz=-3, numpy LU=-3.000000000000003\end{veribox}
\begin{aibox}{AI 讲解（本地 qwen3:8b，仅讲解不产数值）}利用行列式展开或行变换计算三阶行列式的值。\end{aibox}
\vspace{10pt}
\subsection*{题 P05}
\noindent\textbf{题目.}\quad 5. 向量的点积与叉积
已知向量 $\vec{a}=(2,-1,3)$ 与 $\vec{b}=(1,4,-2)$，求二者的点积与叉积，并判断这两个向量是否垂直。\par
\vspace{4pt}\noindent\textbf{解答.}\quad
\begin{enumerate}[leftmargin=1.6em,itemsep=2pt]
  \item 点积：$\vec{u}\cdot\vec{v}=\sum_i u_i v_i$。
  \item 叉积：按三阶行列式展开 $\vec{u}\times\vec{v}$。
  \item 垂直判定：$\vec{u}\cdot\vec{v}\neq 0 \Rightarrow$ 不垂直。
\end{enumerate}
\[ \boxed{\vec{u}\cdot\vec{v}=-8,\quad \vec{u}\times\vec{v}=\left[\begin{matrix}-10 & 7 & 9\end{matrix}\right],\quad \text{两向量不垂直}} \]
\noindent{\small 方法：分量式点积 / 叉积 + 正交判据}\par
\begin{veribox}{独立核验 \okmark}vector\_recheck：矩阵式=-8, 手写分量求和=-8\\cross\_recheck：Matrix.cross=Matrix([[-10, 7, 9]]), 行列式展开=Matrix([[-10, 7, 9]])\end{veribox}
\begin{aibox}{AI 讲解（本地 qwen3:8b，仅讲解不产数值）}利用点积公式和叉积公式分别计算，再通过点积是否为零判断垂直性。\end{aibox}
\vspace{10pt}
\subsection*{题 P06}
\noindent\textbf{题目.}\quad 6. 求一阶微分方程的通解
求微分方程 $y' - 2y = 0$ 的通解。\par
\vspace{4pt}\noindent\textbf{解答.}\quad
\begin{enumerate}[leftmargin=1.6em,itemsep=2pt]
  \item 判定方程阶数与类型：一阶常微分方程。
  \item 分离变量后两端积分，得到含任意常数的通解。
  \item 回代：把解代入原方程，两端之差化简恒为 $0$（见核验栏）。
\end{enumerate}
\[ \boxed{y{\left(x \right)} = C_{1} e^{2 x}} \]
\noindent{\small 方法：一阶常微分方程 —— \texttt{sympy.dsolve}（分离变量 / 特征方程法）}\par
\begin{veribox}{独立核验 \okmark}ode\_residual：解回代原方程，化简后残差恒为 0\end{veribox}
\begin{aibox}{AI 讲解（本地 qwen3:8b，仅讲解不产数值）}识别为齐次线性微分方程，利用指数函数求通解。\end{aibox}
\vspace{10pt}
\subsection*{题 P07}
\noindent\textbf{题目.}\quad 7. 求二阶微分方程满足初值条件的特解
求微分方程 $y'' - 5y' + 6y = 0,\; y(0)=0,\; y'(0)=1$ 的特解。\par
\vspace{4pt}\noindent\textbf{解答.}\quad
\begin{enumerate}[leftmargin=1.6em,itemsep=2pt]
  \item 判定方程阶数与类型：二阶常微分方程。
  \item 写出特征方程并求特征根，按根的情形写通解。
  \item 代入初值条件定出常数，得到满足初值的特解。
  \item 回代：把解代入原方程，两端之差化简恒为 $0$（见核验栏）。
\end{enumerate}
\[ \boxed{y{\left(x \right)} = \left(e^{x} - 1\right) e^{2 x}} \]
\noindent{\small 方法：二阶常微分方程 —— \texttt{sympy.dsolve}（分离变量 / 特征方程法）}\par
\begin{veribox}{独立核验 \okmark}ode\_residual：解回代原方程，化简后残差恒为 0\end{veribox}
\begin{aibox}{AI 讲解（本地 qwen3:8b，仅讲解不产数值）}求解特征方程得到特征根，结合初值条件求特解。\end{aibox}
\vspace{10pt}
\subsection*{题 P08}
\noindent\textbf{题目.}\quad 8. 匀变速直线运动
一辆小车以初速度 $v_0 = 10 m/s$ 出发，沿直线以加速度 $a = 2 m/s^2$ 匀加速运动，
运动时间 t = 5 s，求末速度与这段时间内的位移。\par
\begin{aibox}{未纳入解答}本机确定性求解模板未覆盖该题（已尝试本地模型规范化兜底）；按「不编造答案」原则留空并如实登记。\end{aibox}

\subsection*{题 P09}
\noindent\textbf{题目.}\quad 9. 牛顿第二定律
质量 m = 5 kg 的物体受到合外力 F = 20 N 的作用，求物体获得的加速度。\par
\vspace{4pt}\noindent\textbf{解答.}\quad
\begin{enumerate}[leftmargin=1.6em,itemsep=2pt]
  \item 由牛顿第二定律 $F=ma$，得 $a=F/m$。
  \item $a=20/5=4\,\mathrm{m/s^2}$。
\end{enumerate}
\[ \boxed{a=4\,\mathrm{m/s^2}} \]
\noindent{\small 方法：牛顿第二定律 F = ma}\par
\begin{veribox}{独立核验 \okmark}two\_route：a = F/m = 4；a = F·(1/m) = 4\\units：1 项量纲登记齐备\end{veribox}
\begin{aibox}{AI 讲解（本地 qwen3:8b，仅讲解不产数值）}根据牛顿第二定律 F = ma 直接求加速度。\end{aibox}
\vspace{10pt}
\subsection*{题 P10}
\noindent\textbf{题目.}\quad 10. 动能定理
质量 m = 3 kg 的物体从高 h = 10 m 处由静止落下，取 $g = 9.8 m/s^2$，
求物体落地时的动能，并用落地速度独立换算一次以作核对。\par
\vspace{4pt}\noindent\textbf{解答.}\quad
\begin{enumerate}[leftmargin=1.6em,itemsep=2pt]
  \item 由动能定理 $W_{\text{合}}=\Delta E_k$，重力做功全部转化为动能。
  \item $E_k=mgh=3\times9.8\times10=294\,\mathrm{J}$。
  \item 等价校验：落地速度 $v=\sqrt{2gh}=14\,\mathrm{m/s}$，$E_k=\tfrac12 m v^2=294\,\mathrm{J}$，两路一致。
\end{enumerate}
\[ \boxed{E_k=294\,\mathrm{J}} \]
\noindent{\small 方法：动能定理（重力的功 → 动能增量）}\par
\begin{veribox}{独立核验 \okmark}two\_route：Ek = mgh = 294；Ek = 0.5·m·(2gh) = 294\\units：1 项量纲登记齐备\end{veribox}
\begin{aibox}{AI 讲解（本地 qwen3:8b，仅讲解不产数值）}利用重力势能转化为动能，再通过速度公式核对动能。\end{aibox}
\vspace{10pt}
\subsection*{题 P11}
\noindent\textbf{题目.}\quad 11. 抛体运动
以初速度 $v_0 = 25 m/s$ 斜向上抛出一小球，抛射角 $\theta = 37^\circ$，
取 $g = 9.8 m/s^2$，求射程与最大高度。\par
\vspace{4pt}\noindent\textbf{解答.}\quad
\begin{enumerate}[leftmargin=1.6em,itemsep=2pt]
  \item 水平、竖直分解：$v_{0x}=v_0\cos\theta$，$v_{0y}=v_0\sin\theta$。
  \item 飞行时间 $t=\dfrac{2v_0\sin\theta}{g}=3.0705\,\mathrm{s}$。
  \item 射程 $R=v_0\cos\theta\cdot t=61.305\,\mathrm{m}$。
  \item 最大高度 $H=\dfrac{v_0^2\sin^2\theta}{2g}=11.5491\,\mathrm{m}$。
\end{enumerate}
\[ \boxed{R=61.305\,\mathrm{m},\quad H=11.5491\,\mathrm{m}} \]
\noindent{\small 方法：运动分解（水平匀速 + 竖直匀变速）}\par
\begin{veribox}{独立核验 \okmark}two\_route：R = v0\textasciicircum{}2 sin(2θ)/g = 61.3049551；R = v0 cosθ · (2v0 sinθ/g) = 61.3049551\\units：2 项量纲登记齐备\end{veribox}
\begin{aibox}{AI 讲解（本地 qwen3:8b，仅讲解不产数值）}分解初速度为水平和竖直分量，分别计算射程和最大高度。\end{aibox}
\vspace{10pt}
\subsection*{题 P12}
\noindent\textbf{题目.}\quad 12. 简谐运动的周期
弹簧振子的劲度系数 k = 100 N/m，振子质量 m = 0.25 kg，求该振子的振动周期。\par
\vspace{4pt}\noindent\textbf{解答.}\quad
\begin{enumerate}[leftmargin=1.6em,itemsep=2pt]
  \item 弹簧振子回复力 $F=-kx$，与 $F=ma$ 联立得 $\ddot x+\dfrac km x=0$，故 $\omega=\sqrt{k/m}$。
  \item $\omega=20\,\mathrm{rad/s}\Rightarrow T=\dfrac{2\pi}{\omega}=0.3142\,\mathrm{s}$。
\end{enumerate}
\[ \boxed{T=0.3142\,\mathrm{s}} \]
\noindent{\small 方法：简谐运动角频率（由运动方程推得）}\par
\begin{veribox}{独立核验 \okmark}two\_route：T = 2π/ω, ω=sqrt(k/m) = 0.3141592654；T = 2*pi*sqrt(m/k) = 0.3141592654\\units：1 项量纲登记齐备\end{veribox}
\begin{aibox}{AI 讲解（本地 qwen3:8b，仅讲解不产数值）}根据简谐运动周期公式 T = 2π√(m/k) 计算周期。\end{aibox}
\vspace{10pt}
\subsection*{题 P13}
\noindent\textbf{题目.}\quad 13. 证明题（超范围）
设 $A$ 为可逆矩阵，证明：$\det(A^{-1}) = \dfrac{1}{\det(A)}$。\par
\begin{aibox}{未纳入解答}本机确定性求解模板未覆盖该题（已尝试本地模型规范化兜底）；按「不编造答案」原则留空并如实登记。\end{aibox}

\section*{二、核验汇总}
\begin{center}\small
\begin{tabular}{p{0.5cm}p{6.2cm}p{3.4cm}p{3.0cm}}
\toprule
\# & 题目 & 结果 & 核验 \\
\midrule
P01 & 1. 求解下列线性方程组 求解下列三元线性方程组： … & $x = 1,\quad y = 2,\quad z = 3$ & \textcolor{okgreen}{\textbf{通过}} \\
P02 & 2. 求矩阵的特征值与特征向量 求矩阵 $A = 4 , 1 ; 2 , 3$ 的特征值与对应的特征向量。 & $\lambda_{1}=2,\quad \lambda_{2}=5;\quad \text{对应特征向量见下}$ & \textcolor{okgreen}{\textbf{通过}} \\
P03 & 3. 求矩阵的秩 求矩阵 $A = 1 , 2 , 3 ; 2 , 4 , 6 ; 3 , 6 , 10$ 的秩，并说明行向量组是否线性相… & $r(A)=2$ & \textcolor{okgreen}{\textbf{通过}} \\
P04 & 4. 计算行列式的值 计算下列三阶行列式的值：… & $\det A=-3$ & \textcolor{okgreen}{\textbf{通过}} \\
P05 & 5. 向量的点积与叉积 已知向量 $\vec{a}=(2,-1,3)$ 与 $\vec{b}=(1,4,-2)$，求二者的点积与叉积，并判… & $\vec{u}\cdot\vec{v}=-8,\quad \vec{u}\times\vec{v}=\left[\begin{matrix}-10 & 7 & 9\end{matrix}\right],\quad \text{两向量不垂直}$ & \textcolor{okgreen}{\textbf{通过}} \\
P06 & 6. 求一阶微分方程的通解 求微分方程 $y' - 2y = 0$ 的通解。 & $y{\left(x \right)} = C_{1} e^{2 x}$ & \textcolor{okgreen}{\textbf{通过}} \\
P07 & 7. 求二阶微分方程满足初值条件的特解 求微分方程 $y'' - 5y' + 6y = 0,\; y(0)=0,\; y'(0)=1$ 的… & $y{\left(x \right)} = \left(e^{x} - 1\right) e^{2 x}$ & \textcolor{okgreen}{\textbf{通过}} \\
P08 & 8. 匀变速直线运动 一辆小车以初速度 $v_0 = 10 m/s$ 出发，沿直线以加速度 $a = 2 m/s^2$ 匀加速运动， 运动… & $—$ & 超出范围 \\
P09 & 9. 牛顿第二定律 质量 m = 5 kg 的物体受到合外力 F = 20 N 的作用，求物体获得的加速度。 & $a=4\,\mathrm{m/s^2}$ & \textcolor{okgreen}{\textbf{通过}} \\
P10 & 10. 动能定理 质量 m = 3 kg 的物体从高 h = 10 m 处由静止落下，取 $g = 9.8 m/s^2$， 求物体落地时的… & $E_k=294\,\mathrm{J}$ & \textcolor{okgreen}{\textbf{通过}} \\
P11 & 11. 抛体运动 以初速度 $v_0 = 25 m/s$ 斜向上抛出一小球，抛射角 $\theta = 37^\circ$， 取 … & $R=61.305\,\mathrm{m},\quad H=11.5491\,\mathrm{m}$ & \textcolor{okgreen}{\textbf{通过}} \\
P12 & 12. 简谐运动的周期 弹簧振子的劲度系数 k = 100 N/m，振子质量 m = 0.25 kg，求该振子的振动周期。 & $T=0.3142\,\mathrm{s}$ & \textcolor{okgreen}{\textbf{通过}} \\
P13 & 13. 证明题（超范围） 设 $A$ 为可逆矩阵，证明：$\det(A^{-1}) = \dfrac{1}{\det(A)}$。 & $—$ & 超出范围 \\
\bottomrule
\end{tabular}
\end{center}
\vspace{6pt}
\noindent{\small 本报告由 C4C 作业求解管线自动生成：解答来自确定性符号引擎，核验由独立算法路径二次验算，二者分离；本地大模型仅参与路由复核、题目规范化与讲解生成，不产出最终数值。}\par
\end{document}